% Compile from docs/thesis with latexmk -pdf pptphase1.tex
\documentclass[aspectratio=169,11pt]{beamer}
\usepackage[T1]{fontenc}
\usepackage{lmodern,amsmath,amssymb,mathtools}
\usepackage{graphicx,booktabs,array}
\setbeamertemplate{navigation symbols}{}
\setbeamertemplate{footline}[frame number]
\setbeamercolor{frametitle}{fg=black}
\setbeamerfont{frametitle}{series=\bfseries}
\setbeamersize{text margin left=8mm,text margin right=8mm}
\newcommand{\R}{\mathbb{R}}
\newcommand{\X}{\mathcal{X}}

\title{Learning from Comparison Data}
\subtitle{M.Tech Project: Phase I}
\author{Nitesh Kumar Shah\\ID25M806}
\institute{Indian Institute of Technology Madras\\
    Prof. Suryanarayana Sankagiri}
\date{November 2026}
\titlegraphic{\includegraphics[height=1cm]{images/iitm-logo.png}}


\definecolor{vivateal}{HTML}{12675C}
\setbeamercolor{frametitle}{fg=vivateal}
\setbeamercolor{structure}{fg=vivateal}
\graphicspath{{images/experiments/bradley_terry/}{images/experiments/bradley_terry/optimization/fixed_budget/}{images/experiments/phase1_cdm_six_figures/}}
\newcommand{\vivaplot}[1]{\centering\includegraphics[width=\linewidth,height=.48\textheight,keepaspectratio]{#1}\par}
\newcommand{\takeaway}[1]{\medskip\par\raggedright\small\textbf{#1}}
\begin{document}
\begin{frame}[plain]
    \titlepage
\end{frame}

\begin{frame}{Choice Sets and IIA}
    \textbf{Travel from IIT Madras to the airport}
    \medskip
    \begin{columns}[T,onlytextwidth]
        \column{.48\textwidth}
        \textbf{Original choice set}\\
        Car, taxi, bike, bus
        \column{.48\textwidth}
        \textbf{Expanded choice set}\\
        Car, taxi, bike, bus, walking
    \end{columns}
    \bigskip
    \textbf{Independence of Irrelevant Alternatives (IIA)}
    \medskip

    Under IIA, the relative odds of choosing car over taxi stay the same when walking is added.
    A context effect occurs when the available alternatives change those relative odds.
    \medskip

    {\footnotesize Seshadri et al.\ (2019), \emph{Discovering Context Effects from Raw Choice Data}.}
\end{frame}

\begin{frame}{Bradley--Terry Estimation}
    For items \(i\) and \(j\), the probability that \(i\) defeats \(j\) is
    \[
        P(i\succ j)=\frac{\exp(\theta_i)}{\exp(\theta_i)+\exp(\theta_j)}.
    \]
    \begin{itemize}
        \item Fit the scores by minimising the average negative log-likelihood.
        \item Centre the scores so their sum is zero.
        \item A finite MLE exists when the directed win graph is strongly connected.
    \end{itemize}
\end{frame}

\begin{frame}{Bradley--Terry Parameter Recovery}
    \begin{columns}[T,onlytextwidth]
        \column{.48\textwidth}
        \centering\small\textbf{More comparisons with 10 items}\par
        \vivaplot{recovery_error_vs_comparisons.pdf}
        \column{.48\textwidth}
        \centering\small\textbf{More items at fixed \(m/n\)}\par
        \vivaplot{bt_choix_recovery_vs_items_fixed_ratio.pdf}
    \end{columns}
    \takeaway{RMSE falls as comparisons increase. At fixed \(m/n\), it stays roughly stable as the number of items grows.}
\end{frame}

\begin{frame}{Comparison of Optimisation Methods}
    \begin{columns}[T,onlytextwidth]
        \column{.48\textwidth}
        \centering\small\textbf{Distance from convergence}\par
        \vivaplot{06_distance_from_convergence.pdf}
        \column{.48\textwidth}
        \centering\small\textbf{PyTorch / handcrafted runtime}\par
        \vivaplot{07_pytorch_to_handcrafted_runtime_ratio.pdf}
    \end{columns}
    \takeaway{Recovery errors were similar, but convergence and runtime differed across methods.}
\end{frame}

\begin{frame}{Context-Dependent Model}
    Utility of item \(x\) in choice set \(C\) and its choice probability:
    \[
      u(x\mid C)=\sum_{z\in C\setminus\{x\}}u_{xz},
      \qquad P(x\mid C)=\frac{\exp(u(x\mid C))}{\sum_{y\in C}\exp(u(y\mid C))}.
    \]
    \begin{columns}[T,onlytextwidth]
        \column{.48\textwidth}
        \textbf{Full CDM}\\
        Estimate a parameter for each ordered pair. The centred model has \(n(n-1)-1\) free parameters.
        \column{.48\textwidth}
        \textbf{Low-rank CDM}\\
        Write \(u_{xz}=t_x^\top c_z\) for vectors in \(\R^r\). The implementation stores \(2nr\) trainable entries.
    \end{columns}
\end{frame}

\begin{frame}{Full CDM Parameter Recovery}
    \small Synthetic choices from pairs and triples. The model is fitted on 80\% of the observations.
    \medskip
    \vivaplot{01_full_cdm_recovery.pdf}
    \takeaway{More observations improve recovery for five items. With 50,000 observations fixed, recovery worsens as the number of items increases.}
\end{frame}

\begin{frame}{BTL/Luce Fitted to CDM-Generated Choices}
    \small Prediction gap \(=\) fitted BTL test NLL \(-\) generating CDM test NLL.
    \medskip
    \vivaplot{02_cdm_to_btl_misspecification.pdf}
    \takeaway{More observations do not remove the gap in these experiments. The gap increases with the CDM parameter scale.}
\end{frame}

\begin{frame}{Computational Scaling of Full and Low-Rank CDM}
    \small 250--5,000 items. Each fit uses 1,000 observations and 1,000 Adam updates on the CPU.
    \medskip
    \vivaplot{04_computational_scaling.pdf}
    \raggedright\footnotesize RSS is resident set size, the RAM used by the running process.
    \takeaway{Low-rank CDM stores fewer parameters, but fitting time remains similar in this implementation.}
\end{frame}

\begin{frame}{Unified Comparison on Held-Out Choices}
    \small 10 items, 10,000 observations per generator, 70/15/15 split. Lower test NLL is better.
    \medskip
    \vivaplot{03_unified_model_comparison.pdf}
    \takeaway{The models perform similarly on BTL-generated choices. Both CDM fits have lower test NLL on CDM-generated choices, although the selected rank-five fit has not converged.}
\end{frame}

\begin{frame}{Six-Item CDM Parameter Recovery}
    \small Choice sets contain two to six items. Results average ten runs at each observation count.
    \medskip
    \vivaplot{05_paper_synthetic_recovery.pdf}
    \takeaway{The sum of squared parameter errors decreases as observations increase, while fitting time rises.}
    \medskip

    {\footnotesize Following Seshadri et al.\ (2019).}
\end{frame}

\begin{frame}{Conclusion}
    \textbf{Parameter recovery}\\
    More observations improved recovery. Increasing the number of items under a fixed data budget made estimation more difficult.
    \medskip

    \textbf{Context effects}\\
    Fitting BTL/Luce to the tested CDM-generated choices left a predictive gap.
    \medskip

    \textbf{Low-rank CDM}\\
    Fewer stored parameters did not translate into the same reduction in runtime or total process memory.
    \medskip

    \textbf{Further work}\\
    Check the comparisons across independent datasets and improve optimisation.
\end{frame}
\end{document}
